
%\needspace{4\baselineskip}
\color{unidarkgreen}{\subsection{Блокировка при сквозных токах через ошиновку (БСТО)}}
\color{black}

\begin{enumerate}[label=\arabic{section}.\arabic{subsection}.\arabic{enumi}, labelsep=4pt, leftmargin=0pt, itemindent=57pt, itemsep=0pt, parsep=0pt] % если что можно крутить itemsep 0 или 5
%%%%%%%%%%%%%%%%%%%%%%%%%%%%%%%%%%%%%%%%%%%%%%%%%%%%%%%%%%%%%%%%%%%%%%  ОБЩИЕ СВЕДЕНИЯ %%%%%%%%%%%%%%%%%%%%%%%%%%%%%%%%%%%%%%%%%%%%%%%%%%

\item
Блокировка при сквозных токах через ошиновку (БСТО) предназначена для предотвращения ложного срабатывания быстродействующих защит при протекании сквозных токов внешних КЗ через ошиновку присоединения.

\item
В схемах с двумя выключателями на присоединение при протекании сквозного тока внешнего КЗ через трансформаторы тока двух выключателей может происходить их неодновременное насыщение. Это приводит к искажению формы вторичных токов и возникновению значительного небаланса, который может привести к неправильной работе защит. Функция вводится в работу программным ключом \textbf{«БСТО»}.

\item
Принцип действия БСТО основывается на сравнении фаз между токами первого и второго выключателей. Если токи обоих выключателей превышают уставку \textbf{«Iср БСТО»} и разница фаз токов находится в диапазоне от 90 до 270 градусов в течение времени, задаваемое уставкой \textbf{«Тфикс внеш. КЗ»}, то фиксируется внешнее КЗ и срабатывает блокировка. Выдержка времени \textbf{«Тпрод БСТО»} определяет дополнительную задержку на возврат сигнала фиксации внешнего КЗ. Сигнал срабатывания блокировки действует ограниченное время, определяемое уставкой \textbf{«Тогр БСТО»}. Функция БСТО выполнена пофазной.

\item
Сигнал срабатывания блокировки сбрасывается посредством входных сигналов отключенного положения выключателей.

\item
Блокировка вводится отдельными программными ключами для следующих защит: ДЗ-ФФ-1 и ДЗ-ФЗ-1, ТНЗНП-1, оперативное ускорение ДЗ и ТНЗНП, оперативное ускорение ДЗ и ТНЗНП без выдержки времени, МФТО, ТЗО.

\item
БСТО может быть оперативно выведена из работы от внешнего сигнала \textit{«Вывод БСТО»}, назначенного на соответствующий дискретный вход или функциональную клавишу.

\item
Функциональная схема логики БСТО приведена на рисунке \ref{pict_bk:BSTO}. Уставки БСТО представлены в таблице \ref{app_set:bsto}.

\vspace{2mm}
\begin{figure}[H]
\centering
\includegraphics[width=1\textwidth,height=1\textheight,keepaspectratio]{\fbpath/03.01 ЛВ/561.1481 БСТО/img/BSTO.pdf}
\captionof{figure}{ Функциональная схема БСТО}\label{pict_bk:BSTO}
\end{figure}

\end{enumerate}